%! TEX program = pdflatex
% =====================================================================
%  Short CV / Resume (industry) -- Tyson Lee Swetnam, Ph.D.
%  Updated September 2026. Target length: 2 pages.
%  Overleaf: upload this file, set Compiler to pdfLaTeX.
% =====================================================================
\documentclass[11pt,letterpaper]{article}

\usepackage[letterpaper,top=0.65in,bottom=0.7in,left=0.75in,right=0.75in,footskip=0.35in]{geometry}
\usepackage[T1]{fontenc}
\usepackage[utf8]{inputenc}
\usepackage{lmodern}          % Computer Modern (Latin Modern)
\usepackage{microtype}
\usepackage{fontawesome5}
\usepackage{xcolor}
\usepackage{enumitem}
\usepackage{tabularx}
\usepackage{titlesec}
\usepackage{fancyhdr}
\usepackage{lastpage}
\usepackage{xurl}
\usepackage{hyperref}
\hypersetup{colorlinks=true, linkcolor=black, urlcolor=black, citecolor=black,
  pdftitle={Tyson L. Swetnam -- Resume}, pdfauthor={Tyson L. Swetnam}}

\setlength{\parindent}{0pt}
\setlength{\parskip}{0pt}
\frenchspacing

\pagestyle{fancy}
\fancyhf{}
\renewcommand{\headrulewidth}{0pt}
\fancyfoot[L]{\footnotesize Tyson L. Swetnam \,\textbar\, Resume}
\fancyfoot[R]{\footnotesize\thepage\ of \pageref{LastPage}}

\titleformat{\section}{\large\bfseries}{}{0pt}{\MakeUppercase}[{\vspace{-3pt}\titlerule[0.4pt]}]
\titlespacing*{\section}{0pt}{10pt}{5pt}

% \job{Title}{Dates}{Organization}{Location}
\newcommand{\job}[4]{%
  \begin{tabularx}{\linewidth}{@{}X r@{}}
    \textbf{#1} & {\small #2}\\
    \textit{#3} & {\small #4}
  \end{tabularx}\par}
\newcommand{\jobline}[3]{%
  \begin{tabularx}{\linewidth}{@{}X r@{}}
    \textbf{#1}, \textit{#2} & {\small #3}
  \end{tabularx}\par\vspace{1pt}}
\newcommand{\me}{\textbf{Swetnam, T.L.}}
\newcommand{\doi}[1]{\href{https://doi.org/#1}{\nolinkurl{doi:#1}}}

\setlist[itemize]{leftmargin=1.2em,itemsep=1pt,parsep=0pt,topsep=2pt,label={\small\textbullet}}

\begin{document}

% ---------- Header ----------------------------------------------------
\begin{center}
  {\fontsize{22}{26}\selectfont\bfseries Tyson L. Swetnam, \textmd{Ph.D.}}\par\vspace{4pt}
  {\normalsize Research Computing, Data Platforms \& Applied AI Leader}\par\vspace{5pt}
  {\small
   \faMapMarker*\,Albuquerque, NM \quad
   \faEnvelope[regular]\,\href{mailto:tswetnam@unm.edu}{tswetnam@unm.edu} \quad
   \faGlobe\,\href{https://tysonswetnam.com}{tysonswetnam.com}\par\vspace{2pt}
   \faGithub\,\href{https://github.com/tyson-swetnam}{tyson-swetnam} \quad
   \faLinkedin\,\href{https://www.linkedin.com/in/tyson-swetnam}{tyson-swetnam} \quad
   \faGraduationCap\,\href{https://scholar.google.com/citations?user=HeftZR8AAAAJ}{Scholar}}
\end{center}

% ---------- Summary ---------------------------------------------------
\section{Summary}
Technical leader with 15+ years building and operating cloud, HPC, and data infrastructure for scientific research. Director of the University of New Mexico's Center for Advanced Research Computing (CARC). Former Co-PI of CyVerse, the NSF-funded national cyberinfrastructure for the life sciences. Track record of winning and delivering large federal programs, standing up platforms that serve thousands of researchers, and translating emerging technology (containers, cloud-native data, LLMs, and agentic AI) into production services and hands-on training.

% ---------- Skills ----------------------------------------------------
\section{Core Competencies}
\begin{tabularx}{\linewidth}{@{}p{1.55in}X@{}}
\textbf{Leadership} & Research-computing strategy \& operations; program management of multi-institution federal awards; budgets \& proposals; team building; stakeholder \& funder relations (NSF, USDA, NIH, NASA) \\[2pt]
\textbf{AI \& ML} & Large language models, agentic AI \& automation, metadata-aware AI agents, knowledge graphs \& ontologies, GPU computing, prompt engineering \\[2pt]
\textbf{Data platforms} & Federated data mesh, lakehouse table formats (Apache Iceberg, DuckLake), iRODS, FAIR data \& metadata curation, cloud-optimized \& analysis-ready data \\[2pt]
\textbf{Infrastructure} & HPC, public research clouds (Jetstream2, ACCESS), Docker, Singularity/Apptainer, Kubernetes, Linux, GitHub CI/CD \\[2pt]
\textbf{Languages \& tools} & Python, R, Bash, SQL, Jupyter, Google Earth Engine, GIS \\[2pt]
\textbf{Geospatial} & Satellite, airborne \& drone (sUAS) remote sensing; lidar \& photogrammetry; environmental monitoring
\end{tabularx}

% ---------- Experience ------------------------------------------------
\section{Experience}

\job{Director, Center for Advanced Research Computing}{7/2026 -- Present}{University of New Mexico \textbullet{} also Associate Professor of Computer Science}{Albuquerque, NM}
\begin{itemize}
  \item Lead UNM's central HPC, storage, and research-computing center, which serves faculty and students across the state's flagship university.
  \item Setting direction for open-access research data storage and for making advanced computing and AI accessible to all New Mexico researchers.
  \item PI of \textbf{NSF MESA} (\$4.6M, 2026--2029), a NAIRR research-infrastructure platform in which metadata-aware AI agents describe, organize, and connect scientific data across fields.
\end{itemize}
\vspace{4pt}

\job{Research Associate Professor}{7/2023 -- 6/2026}{University of Arizona, Institute for Artificial Intelligence \& Society}{Tucson, AZ}
\begin{itemize}
  \item Led the University of Arizona cyberinfrastructure work for two NSF synthesis centers: ESIIL (\$20M program; \$1.46M UA scope) and NCEMS (\$1.93M UA scope).
  \item Senior personnel on the USDA--NSF AI Institute for Resilient Agriculture (AIIRA, \$20M) and on the Data Management Core of the NIH Superfund DUST Center (\$14.8M).
  \item Created and taught LLM and agentic-AI workshops and course materials (GPT-101; AI Automation \& Agents) for researchers, staff, and students.
\end{itemize}
\vspace{4pt}

\job{Research Assistant Professor \& Data Scientist, CyVerse}{9/2016 -- 6/2023}{University of Arizona, BIO5 Institute}{Tucson, AZ}
\begin{itemize}
  \item Co-PI of \textbf{CyVerse} (NSF, \$15.2M), national cyberinfrastructure providing data storage, cloud analysis, and containerized apps to the life-science research community.
  \item Built and ran Container Camp and Foundational Open Science Skills training, teaching Docker, Kubernetes, Singularity, and reproducible workflows to researchers nationwide.
  \item Delivered scalable GPU and ML pipelines for crop phenomics (PhytoOracle) and drone data, and secured TACC Frontera and Longhorn GPU allocations.
  \item Senior personnel on the NSF HydroGEN (\$5M) and Groundwater Data Platform (\$1M) machine-learning water-resource projects.
\end{itemize}
\vspace{4pt}

\jobline{Associate Research Scientist / Postdoctoral Researcher}{University of Arizona \& University of Utah}{2014 -- 2016}
\jobline{Fire Management Specialist}{USDA Forest Service, Coronado National Forest}{2008 -- 2012}
\jobline{Forestry Technician (Wildland Firefighter)}{U.S.\ Dept.\ of the Interior, Saguaro National Park}{2002 -- 2005}

% ---------- Education -------------------------------------------------
\section{Education}
\jobline{Ph.D., Natural Resources (minor: Remote Sensing \& Spatial Analysis)}{University of Arizona}{2013}
\jobline{M.S., Natural Resources (GIS Technical Certificate)}{University of Arizona}{2006}
\jobline{B.S., Ecology \& Evolutionary Biology}{University of Arizona}{2002}

% ---------- Impact ----------------------------------------------------
\section{Selected Impact \& Recognition}
\begin{itemize}
  \item PI, Co-PI, or senior personnel on 25+ funded projects, including four programs of \$14M--\$20M each.
  \item 40 peer-reviewed publications; 2,000+ citations (h-index 26, Google Scholar).
  \item Best Short Paper, Workforce Development, PEARC '23; ABOR TRIF Eighteenth Mile Award (\$141K, 2023) for AEOLUS, a drone-data laboratory.
  \item Grant review panelist for NSF CISE, NSF DBI, and USDA NIFA; member of the NSF CI Compass Cloud Infrastructure Working Group; certified Carpentries instructor.
  \item Invited speaker at TACC, NASA JPL, NVIDIA GPU Cloud, the Santa Fe Institute, the University of Cambridge, and NSF--CNPq (Brazil).
\end{itemize}

% ---------- Publications ----------------------------------------------
\section{Selected Publications}
\begin{itemize}[label={}, leftmargin=1.2em, itemindent=-1.2em, itemsep=2pt]
  \item \me, P.B. Antin, R. Bartelme, et al.\ (2024) CyVerse: Cyberinfrastructure for open science. \textit{PLOS Computational Biology}. \doi{10.1371/journal.pcbi.1011270}
  \item Roy, A., E. Ward, I. Choi, et al.\ (2025) MDRepo -- an open data warehouse for community-contributed molecular dynamics simulations of proteins. \textit{Nucleic Acids Research} 53(D1). \doi{10.1093/nar/gkae1109}
  \item Thessen, A., L. Cooper, \me, et al.\ (2023) Using knowledge graphs to infer gene expression in plants. \textit{Frontiers in Artificial Intelligence} 6:92. \doi{10.3389/frai.2023.1201002}
  \item Skidmore, E., M. Cosi, \me, et al.\ (2023) Cloud computing for research and education gets a sweet upgrade with CACAO. \textit{PEARC '23}. ACM. \doi{10.1145/3569951.3597555}
  \item Gonzalez, E.M., A. Zarei, N. Hendler, et al.\ (2023) PhytoOracle: Scalable, modular phenomics data processing pipelines. \textit{Frontiers in Plant Science} 14. \doi{10.1002/essoar.10508789.1}
\end{itemize}

\vfill
{\centering\footnotesize Full CV and publication list at \href{https://tysonswetnam.com/cv/}{tysonswetnam.com/cv}\par}

\end{document}
